\documentclass[11pt,a4paper]{article}
\usepackage{fontspec,amsmath,amssymb,amsthm,microtype}
\setmainfont{lmroman10-regular.otf}[BoldFont=lmroman10-bold.otf,ItalicFont=lmroman10-italic.otf,BoldItalicFont=lmroman10-bolditalic.otf]
\usepackage{xcolor}
\usepackage[most]{tcolorbox}
\usepackage{mathrsfs}
\providecommand{\lt}{<}
\providecommand{\gt}{>}
\setmonofont{DejaVuSansMono.ttf}[BoldFont=DejaVuSansMono-Bold.ttf,ItalicFont=DejaVuSansMono-Oblique.ttf,BoldItalicFont=DejaVuSansMono-BoldOblique.ttf,Scale=0.86]
\newfontfamily\leanfallback{STIXTwoMath-Regular.otf}[Scale=0.86]
\definecolor{leanrule}{gray}{0.62}
\definecolor{leanback}{gray}{0.962}
\definecolor{leantitle}{gray}{0.22}
\definecolor{leanmuted}{gray}{0.38}
\newtcolorbox{leanbox}[1]{enhanced,breakable,colback=leanback,colframe=leanrule,
  boxrule=0.4pt,arc=0.6mm,left=2.4mm,right=2.4mm,top=1.8mm,bottom=1.8mm,
  before skip=0.8em,after skip=1em,fonttitle=\small\bfseries,coltitle=leantitle,
  colbacktitle=leanback,titlerule=0.3pt,title={#1}}
\newlength{\leancharwidth}
\newenvironment{leancode}{\par\smallskip\begingroup\ttfamily\footnotesize\raggedright
  \settowidth{\leancharwidth}{M}\setlength{\parskip}{0pt}\setlength{\parindent}{0pt}
  \frenchspacing}{\par\endgroup\smallskip}
\newcommand{\leanline}[2]{\hangindent=\dimexpr#1\leancharwidth+4\leancharwidth\relax
  \hangafter=1\noindent\hspace*{#1\leancharwidth}#2\par}
\newcommand{\leansource}[1]{\par{\footnotesize\color{leanmuted}#1}\par}
\newcommand{\leanlabel}[1]{\par\smallskip\noindent{\small\bfseries #1}\par}
\usepackage[margin=25mm]{geometry}
\usepackage[unicode,hidelinks]{hyperref}
\hypersetup{pdfauthor={LangevinLean Project}}
\setlength{\emergencystretch}{3em}
\setlength{\parskip}{0.4em}
\setlength{\parindent}{0pt}
\allowdisplaybreaks[2]
\raggedbottom
\title{\texorpdfstring{Langevin Monte Carlo with Average Smoothness}{Langevin Monte Carlo with Average Smoothness}}
\author{LangevinLean Project}
\date{}
\begin{document}\maketitle
Draft mathematical exposition; these notes may contain mistakes.
The main sampling theorems and decreasing-step proposition are proved for actual Gaussian laws with the source hypotheses and constants. The source audit records broader auxiliary coverage separately.
\section{\texorpdfstring{Theorem}{Theorem}}


Let \(f:\mathbb R^p\to\mathbb R\) be differentiable,
\(m\)-strongly convex with \(M_\infty\)-Lipschitz gradient, where
\(m>0\). The coordinate constants are the literal Taylor suprema in the
paper, and \(M_{\rm av}=p^{-1}\sum_j M_j\). For the actual independent
Gaussian step \(X^+=X-h\nabla f(X)+\sqrt{2h}\,Z\), every \(h\ge0\) satisfies


\[W_2^2(\nu^+,\pi)\le
\frac{1-mh}{1+mh}W_2^2(\nu,\pi)+
\frac{(M_{\rm av}+m)h^2p}{1+mh}+
\frac{h^2(hM_\infty-1)}{1+mh}\mathbb E_\nu\|\nabla f\|^2.\]


The energy comparison, Gaussian smoothing, Brownian orthogonality and
transport infimum are proved from these hypotheses. This real-valued local
formula is stated for finite quadratic transport, so every signed energy term
is defined. The nonnegative global estimates below also cover infinite
initial transport.

\begin{leanbox}{Lean: \texttt{lmc\_\allowbreak{}one\_\allowbreak{}step}}
\leansource{\texttt{lmc\_\allowbreak{}one\_\allowbreak{}step} \quad AverageLMC/\allowbreak{}GaussianLMC.lean, lines 19–50}
\begin{leancode}
\leanline{0}{theorem\ lmc\_one\_step\ \{p\ :\ ℕ\}\ (hp\ :\ 0\ <\ p)\ (f\ :\ State\ p\ →\ ℝ)}
\leanline{4}{(m\ h\ :\ ℝ)\ (L\ :\ ℝ≥0)\ (hm\ :\ 0\ <\ m)\ (hh\ :\ 0\ ≤\ h)}
\leanline{4}{(hf\ :\ Differentiable\ ℝ\ f)\ (hsc\ :\ StrongConvexOn\ univ\ m\ f)}
\leanline{4}{(hLip\ :\ LipschitzWith\ L\ (gradient\ f))}
\leanline{4}{(μ\ :\ Measure\ (State\ p))\ [IsProbabilityMeasure\ μ]}
\leanline{4}{(hfinite\ :\ wasserstein2\ μ\ (gibbsMeasure\ f)\ ≠\ ⊤)\ :}
\leanline{4}{wasserstein2Real\ (ulaLawStep\ f\ h\ μ)\ (gibbsMeasure\ f)\ \textasciicircum{}\ 2\ ≤}
\leanline{6}{(1\ -\ m\ *\ h)\ /\ (1\ +\ m\ *\ h)\ *\ wasserstein2Real\ μ\ (gibbsMeasure\ f)\ \textasciicircum{}\ 2\ +}
\leanline{6}{((averageSmoothness\ f\ +\ m)\ *\ h\ \textasciicircum{}\ 2\ *\ p)\ /\ (1\ +\ m\ *\ h)\ +}
\leanline{6}{h\ \textasciicircum{}\ 2\ *\ (h\ *\ (L\ :\ ℝ)\ -\ 1)\ /\ (1\ +\ m\ *\ h)\ *}
\leanline{8}{(∫\ x,\ ‖gradient\ f\ x‖\ \textasciicircum{}\ 2\ ∂μ)}
\end{leancode}
\end{leanbox}

\section{\texorpdfstring{Gaussian smoothing and the energy comparison}{Gaussian smoothing and the energy comparison}}



\textbf{Objects and hypotheses.} Let \(p\ge1\), \(m>0\), and let \(f:\mathbb R^p\to\mathbb R\) be differentiable, \(m\)-strongly convex and have an \(M_\infty\)-Lipschitz gradient. Thus

\[f(y)\ge f(x)+\langle\nabla f(x),y-x\rangle+\frac m2\|y-x\|^2,\qquad\|\nabla f(x)-\nabla f(y)\|\le M_\infty\|x-y\|.\]

The Gibbs probability law is \(\pi(dx)=Z^{-1}e^{-f(x)}dx\), where strong convexity makes \(0<Z<\infty\) and gives finite second moments. A coupling of two probability laws is a joint law with those marginals, and \(W_2\) is the infimum of the root-mean-square displacement over couplings. Work first with \(W_2(\nu,\pi)<\infty\). Throughout, a norm of a random vector means \(\|X\|_{L^2}=(\mathbb E\|X\|_2^2)^{1/2}\); its inner product is \(\mathbb E\langle X,Y\rangle\). The gradient bound \(\|\nabla f(x)\|\le\|\nabla f(0)\|+M_\infty\|x\|\) makes all score second moments finite for the states considered. The two Taylor bounds give \(|f(x)|\le|f(0)|+\|\nabla f(0)\|\|x\|+M_\infty\|x\|^2/2\), justifying the expected potential differences.




The
supremum defining \(M_j\) gives
\(f(x+ae_j)\le f(x)+a\,\partial_jf(x)+\tfrac12M_ja^2\) for all \(x\) and
\(a\). Insert the coordinates of a standard Gaussian vector \(\xi\) one at a
time, with \(a=\sqrt{2s}\,\xi_j\). Each \(\xi_j\) is centered and independent
of the coordinates already inserted, so its linear term has mean zero, while
\(\mathbb E a^2=2s\). Summing over \(j\),


\[\mathbb E f(x+\sqrt{2s}\,\xi)-f(x)\le s\sum_{j=1}^pM_j=pM_{\rm av}s
\qquad(s\ge0).\]


Couple \(\vartheta\sim\nu\) with \(Y_0\sim\pi\) at finite cost. Let
\(Y_t=Y_0-\int_0^t\nabla f(Y_s)\,ds+B_t\), \(B_t=\sqrt2W_t\), be the
stationary Langevin diffusion driven by a Brownian motion independent of
\((\vartheta,Y_0)\). Put \(A=\vartheta-h\nabla f(\vartheta)\), so that
\(A+B_h\sim\nu^+\) and \(Y_t\sim\pi\). Write \(w=W_2^2(\nu^+,\pi)\) and use
\(L^2\) norms. The common noise cancels:


\[\|A+B_h-Y_h\|^2=\|A-Y_0\|^2+2\int_0^h\langle A+B_t-Y_t,\nabla f(Y_t)\rangle\,dt.\]


Strong convexity bounds the integrand by
\(2\mathbb E[f(A+B_t)-f(Y_t)]-m\|A+B_t-Y_t\|^2\). Stationarity and smoothing at
the independent centre \(A\) give
\(\mathbb E[f(A+B_t)-f(Y_t)]\le\mathbb E[f(A)-f(Y_0)]+pM_{\rm av}t\). The
future increment \(B_h-B_t\) is orthogonal to \(A+B_t-Y_t\) and has squared
norm \(2(h-t)p\). Since \(w\le\|A+B_h-Y_t\|^2\) for every \(t\), this gives
\(\|A+B_t-Y_t\|^2\ge w-2(h-t)p\). Integrating over \([0,h]\) yields


\[(1+mh)\,w\le\|A-Y_0\|^2+2h\,\mathbb E[f(A)-f(Y_0)]+(M_{\rm av}+m)h^2p.\]\begin{leanbox}{Lean: \texttt{integral\_\allowbreak{}stdGaussian\_\allowbreak{}smoothing}, \texttt{gaussian\_\allowbreak{}energy\_\allowbreak{}bound}}
\leansource{\texttt{integral\_\allowbreak{}stdGaussian\_\allowbreak{}smoothing} \quad AverageLMC/\allowbreak{}Smoothing.lean, lines 154–175}
\begin{leancode}
\leanline{0}{theorem\ integral\_stdGaussian\_smoothing\ \{p\ :\ ℕ\}\ (hp\ :\ 0\ <\ p)}
\leanline{4}{(f\ :\ State\ p\ →\ ℝ)\ (M\ :\ ℝ≥0)\ (hf\ :\ Differentiable\ ℝ\ f)}
\leanline{4}{(hLip\ :\ LipschitzWith\ M\ (gradient\ f))\ (x\ :\ State\ p)\ (s\ :\ ℝ)\ (hs\ :\ 0\ ≤\ s)\ :}
\leanline{4}{(∫\ z,\ f\ (x\ +\ Real.sqrt\ (2\ *\ s)\ •\ z)\ -\ f\ x\ ∂stdGaussian\ (State\ p))\ ≤}
\leanline{6}{averageSmoothness\ f\ *\ (p\ :\ ℝ)\ *\ s}
\end{leancode}
\leansource{\texttt{gaussian\_\allowbreak{}energy\_\allowbreak{}bound} \quad AverageLMC/\allowbreak{}GaussianPerturbation.lean, lines 58–106}
\begin{leancode}
\leanline{0}{theorem\ gaussian\_energy\_bound\ (f\ :\ E\ →\ ℝ)\ (m\ h\ d\ a\ b\ w\ :\ ℝ)\ (L\ :\ ℝ≥0)}
\leanline{4}{(hf\ :\ Differentiable\ ℝ\ f)\ (hsc\ :\ StrongConvexOn\ univ\ m\ f)}
\leanline{4}{(hLip\ :\ LipschitzWith\ L\ (gradient\ f))\ (hm\ :\ 0\ ≤\ m)\ (hh\ :\ 0\ ≤\ h)}
\leanline{4}{(X\ :\ Lp\ E\ 2\ P)\ (Y\ B\ :\ ℝ\ →\ Lp\ E\ 2\ P)\ (hY\ :\ Continuous\ Y)\ (hB\ :\ Continuous\ B)}
\leanline{4}{(hYe\ :\ ∀\ t\ ∈\ Icc\ 0\ h,\ Y\ t\ -\ Y\ 0\ =}
\leanline{6}{-(∫\ s\ in\ 0..t,\ ULA.lpDrift\ P\ (gradient\ f)\ hLip\ (Y\ s))\ +\ B\ t)}
\leanline{4}{(horth\ :\ ∀\ t\ ∈\ Icc\ 0\ h,\ {\leanfallback ⟪}X\ +\ B\ t\ -\ Y\ t,\ B\ h\ -\ B\ t{\leanfallback ⟫}\ =\ 0)}
\leanline{4}{(hnoise\ :\ ∀\ t\ ∈\ Icc\ 0\ h,\ ‖B\ h\ -\ B\ t‖\ \textasciicircum{}\ 2\ =\ 2\ *\ (h\ -\ t)\ *\ d)}
\leanline{4}{(hcost\ :\ ∀\ t\ ∈\ Icc\ 0\ h,\ w\ ≤\ ‖X\ +\ B\ h\ -\ Y\ t‖\ \textasciicircum{}\ 2)}
\leanline{4}{(henergy\ :\ ∀\ t\ ∈\ Icc\ 0\ h,}
\leanline{6}{(∫\ ω,\ f\ ((X\ +\ B\ t)\ ω)\ ∂P)\ -\ (∫\ ω,\ f\ (Y\ t\ ω)\ ∂P)\ ≤\ a\ +\ b\ *\ t)\ :}
\leanline{4}{(1\ +\ m\ *\ h)\ *\ w\ ≤\ ‖X\ -\ Y\ 0‖\ \textasciicircum{}\ 2\ +\ 2\ *\ h\ *\ a\ +\ h\ \textasciicircum{}\ 2\ *\ b\ +\ m\ *\ h\ \textasciicircum{}\ 2\ *\ d}
\end{leancode}
\leanlabel{How the Lean proof follows the argument}
\texttt{integral\_\allowbreak{}std\allowbreak{}Gaussian\_\allowbreak{}smoothing} specializes the independent-coordinate smoothing theorem to the standard Gaussian law. In its finite-coordinate induction, centering removes the linear term and the unit second moment evaluates the quadratic term. \texttt{gaussian\_\allowbreak{}energy\_\allowbreak{}bound} inserts this smoothing inequality into the common-noise stationary comparison; the Brownian orthogonality identity contributes the extra curvature-times-dimension term. All these are integral statements with integrability discharged before linearity is used.
\end{leanbox}

\section{\texorpdfstring{The one-step Wasserstein bound}{The one-step Wasserstein bound}}


With
\(G=\mathbb E\|\nabla f(\vartheta)\|^2\), strong convexity gives the
gradient-step inequality


\[\|A-Y_0\|^2\le(1-mh)\|\vartheta-Y_0\|^2+2h\,\mathbb E[f(Y_0)-f(\vartheta)]+h^2G,\]


and the \(M_\infty\)-Lipschitz gradient gives the descent inequality
\(2\mathbb E[f(A)-f(\vartheta)]\le-h(2-hM_\infty)G\). Adding both to the energy
comparison cancels every potential value:


\[(1+mh)\,W_2^2(\nu^+,\pi)\le(1-mh)\|\vartheta-Y_0\|^2+(M_{\rm av}+m)h^2p
+h^2(hM_\infty-1)G.\]


The output law does not depend on the initial coupling. Taking the
infimum over finite-cost couplings, which is valid also when \(1-mh<0\),
and dividing by \(1+mh\) gives Theorem 1's one-step bound for \(h>0\). At
\(h=0\) the transition is the identity.




\textbf{The cancellation explicitly.} Strong convexity at \(\vartheta\), with comparison point \(Y_0\), gives

\[-2h\mathbb E\langle\nabla f(\vartheta),\vartheta-Y_0\rangle\le2h\mathbb E[f(Y_0)-f(\vartheta)]-mh\|\vartheta-Y_0\|_{L^2}^2.\]

Expanding \(\|\vartheta-h\nabla f(\vartheta)-Y_0\|_{L^2}^2\) produces the gradient-step bound. For the descent bound apply the upper Taylor model at \(\vartheta\) with displacement \(-h\nabla f(\vartheta)\): its linear term is \(-h\|\nabla f(\vartheta)\|^2\) and its quadratic term is \(M_\infty h^2\|\nabla f(\vartheta)\|^2/2\). Adding the two inequalities to the energy comparison cancels \(f(Y_0)\), \(f(\vartheta)\) and \(f(A)\) exactly. The remaining score coefficient is \(h^2-2h^2+h^3M_\infty=h^2(hM_\infty-1)\). Every denominator is positive because \(m>0\) and \(h\ge0\).


\begin{leanbox}{Lean: \texttt{lp\_\allowbreak{}gradient\_\allowbreak{}step}, \texttt{lp\_\allowbreak{}gradient\_\allowbreak{}descent}, \texttt{lmc\_\allowbreak{}one\_\allowbreak{}step}}
\leansource{\texttt{lp\_\allowbreak{}gradient\_\allowbreak{}step} \quad AverageLMC/\allowbreak{}ExpectedGradientStep.lean, lines 15–33}
\begin{leancode}
\leanline{0}{theorem\ lp\_gradient\_step\ (f\ :\ E\ →\ ℝ)\ (m\ h\ :\ ℝ)\ (L\ :\ ℝ≥0)}
\leanline{4}{(hf\ :\ Differentiable\ ℝ\ f)\ (hsc\ :\ StrongConvexOn\ Set.univ\ m\ f)}
\leanline{4}{(hLip\ :\ LipschitzWith\ L\ (gradient\ f))\ (hh\ :\ 0\ ≤\ h)\ (X\ Y\ :\ Lp\ E\ 2\ P)\ :}
\leanline{4}{‖X\ -\ h\ •\ ULA.lpDrift\ P\ (gradient\ f)\ hLip\ X\ -\ Y‖\ \textasciicircum{}\ 2\ ≤}
\leanline{6}{(1\ -\ m\ *\ h)\ *\ ‖X\ -\ Y‖\ \textasciicircum{}\ 2\ +}
\leanline{6}{2\ *\ h\ *\ ((∫\ ω,\ f\ (Y\ ω)\ ∂P)\ -\ ∫\ ω,\ f\ (X\ ω)\ ∂P)\ +}
\leanline{6}{h\ \textasciicircum{}\ 2\ *\ ‖ULA.lpDrift\ P\ (gradient\ f)\ hLip\ X‖\ \textasciicircum{}\ 2}
\end{leancode}
\leansource{\texttt{lp\_\allowbreak{}gradient\_\allowbreak{}descent} \quad AverageLMC/\allowbreak{}ExpectedGradientStep.lean, lines 37–63}
\begin{leancode}
\leanline{0}{theorem\ lp\_gradient\_descent\ (f\ :\ E\ →\ ℝ)\ (h\ :\ ℝ)\ (L\ :\ ℝ≥0)}
\leanline{4}{(hf\ :\ Differentiable\ ℝ\ f)\ (hLip\ :\ LipschitzWith\ L\ (gradient\ f))}
\leanline{4}{(hh\ :\ 0\ ≤\ h)\ (X\ :\ Lp\ E\ 2\ P)\ :}
\leanline{4}{2\ *\ ((∫\ ω,\ f\ ((X\ -\ h\ •\ ULA.lpDrift\ P\ (gradient\ f)\ hLip\ X)\ ω)\ ∂P)\ -}
\leanline{6}{∫\ ω,\ f\ (X\ ω)\ ∂P)\ ≤}
\leanline{6}{-h\ *\ (2\ -\ h\ *\ (L\ :\ ℝ))\ *\ ‖ULA.lpDrift\ P\ (gradient\ f)\ hLip\ X‖\ \textasciicircum{}\ 2}
\end{leancode}
\leansource{\texttt{lmc\_\allowbreak{}one\_\allowbreak{}step} \quad AverageLMC/\allowbreak{}GaussianLMC.lean, lines 19–50}
\begin{leancode}
\leanline{0}{theorem\ lmc\_one\_step\ \{p\ :\ ℕ\}\ (hp\ :\ 0\ <\ p)\ (f\ :\ State\ p\ →\ ℝ)}
\leanline{4}{(m\ h\ :\ ℝ)\ (L\ :\ ℝ≥0)\ (hm\ :\ 0\ <\ m)\ (hh\ :\ 0\ ≤\ h)}
\leanline{4}{(hf\ :\ Differentiable\ ℝ\ f)\ (hsc\ :\ StrongConvexOn\ univ\ m\ f)}
\leanline{4}{(hLip\ :\ LipschitzWith\ L\ (gradient\ f))}
\leanline{4}{(μ\ :\ Measure\ (State\ p))\ [IsProbabilityMeasure\ μ]}
\leanline{4}{(hfinite\ :\ wasserstein2\ μ\ (gibbsMeasure\ f)\ ≠\ ⊤)\ :}
\leanline{4}{wasserstein2Real\ (ulaLawStep\ f\ h\ μ)\ (gibbsMeasure\ f)\ \textasciicircum{}\ 2\ ≤}
\leanline{6}{(1\ -\ m\ *\ h)\ /\ (1\ +\ m\ *\ h)\ *\ wasserstein2Real\ μ\ (gibbsMeasure\ f)\ \textasciicircum{}\ 2\ +}
\leanline{6}{((averageSmoothness\ f\ +\ m)\ *\ h\ \textasciicircum{}\ 2\ *\ p)\ /\ (1\ +\ m\ *\ h)\ +}
\leanline{6}{h\ \textasciicircum{}\ 2\ *\ (h\ *\ (L\ :\ ℝ)\ -\ 1)\ /\ (1\ +\ m\ *\ h)\ *}
\leanline{8}{(∫\ x,\ ‖gradient\ f\ x‖\ \textasciicircum{}\ 2\ ∂μ)}
\end{leancode}
\leanlabel{How the Lean proof follows the argument}
\texttt{lp\_\allowbreak{}gradient\_\allowbreak{}step} and \texttt{lp\_\allowbreak{}gradient\_\allowbreak{}descent} are the two expected deterministic inequalities whose potential terms cancel. \texttt{lmc\_\allowbreak{}one\_\allowbreak{}step} first separates the zero step. In the positive branch, \texttt{hsm} is the exact coordinate smoothing estimate, \texttt{houtput} proves the output transport finite, and \texttt{wasserstein2\_\allowbreak{}sq\_\allowbreak{}bound\_\allowbreak{}of\_\allowbreak{}finite\_\allowbreak{}couplings} passes the common bound to the infimum even when its input-cost coefficient has either sign. The final field normalization divides by the strictly positive denominator.
\end{leanbox}

\section{\texorpdfstring{Constant-step convergence}{Constant-step convergence}}


For
\(0\le h\le1/M_\infty\) the score term is nonpositive and is dropped. The
coordinate bounds give \(m\le M_{\rm av}\le M_\infty\), so \(t=mh\in[0,1]\).
Finite second moments propagate from \(\nu_0\), so
\(r_k=W_2^2(\nu_k,\pi)\) satisfies, at every step,


\[r_{k+1}\le q\,r_k+\frac{(M_{\rm av}+m)h^2p}{1+t},\qquad q=\frac{1-t}{1+t}.\]


Since \(1-q=2t/(1+t)\), the fixed point is
\(c=(M_{\rm av}+m)hp/(2m)\), and induction gives
\(r_k\le q^kr_0+(1-q^k)c\le q^kr_0+c\). Finally \(q\le e^{-2t}\) on
\([0,1]\), so


\[W_2^2(\nu_k,\pi)\le e^{-2mkh}W_2^2(\nu_0,\pi)+\frac{(M_{\rm av}+m)hp}{2m}.\]


If \(W_2(\nu_0,\pi)=\infty\), the same bound holds trivially in
\([0,\infty]\).




\textbf{Induction explicitly.} If \(h=0\), the chain is constant and the assertion follows by equality. If \(h>0\), write \(F=(M_{\rm av}+m)h^2p/(1+mh)\). Since \(0\le q<1\), the identity \(F=(1-q)c\) holds with \(c=(M_{\rm av}+m)hp/(2m)\). At \(k=0\) the proposed sharper bound \(r_k\le q^kr_0+(1-q^k)c\) is equality. If it holds at \(k\), multiply it by \(q\ge0\), add \((1-q)c\), and obtain \(q^{k+1}r_0+(1-q^{k+1})c\). This proves the induction. For \(0\le t<1\), differentiation gives

\[\frac{d}{dt}\left(\log\frac{1-t}{1+t}+2t\right)=-\frac2{1-t^2}+2=-\frac{2t^2}{1-t^2}\le0.\]

The expression vanishes at zero, so \(q\le e^{-2t}\). At \(t=1\), \(q=0\) and the same inequality is immediate. Raising this nonnegative comparison to the natural power \(k\) yields the displayed exponential rate.


\begin{leanbox}{Lean: \texttt{contraction\_\allowbreak{}le\_\allowbreak{}exp}, \texttt{constant\_\allowbreak{}step\_\allowbreak{}normalized}, \texttt{lmc\_\allowbreak{}constant\_\allowbreak{}step}}
\leansource{\texttt{contraction\_\allowbreak{}le\_\allowbreak{}exp} \quad AverageLMC/\allowbreak{}Recurrence.lean, lines 21–34}
\begin{leancode}
\leanline{0}{theorem\ contraction\_le\_exp\ \{t\ :\ ℝ\}\ (ht\ :\ 0\ ≤\ t)\ (ht1\ :\ t\ ≤\ 1)\ :}
\leanline{4}{(1\ -\ t)\ /\ (1\ +\ t)\ ≤\ Real.exp\ (-2\ *\ t)}
\end{leancode}
\leansource{\texttt{constant\_\allowbreak{}step\_\allowbreak{}normalized} \quad AverageLMC/\allowbreak{}Recurrence.lean, lines 78–83}
\begin{leancode}
\leanline{0}{theorem\ constant\_step\_normalized\ (r\ :\ ℕ\ →\ ℝ)\ (a\ t\ :\ ℝ)}
\leanline{4}{(ha\ :\ 0\ ≤\ a)\ (ht\ :\ 0\ ≤\ t)\ (ht1\ :\ t\ ≤\ 1)\ (hr0\ :\ 0\ ≤\ r\ 0)}
\leanline{4}{(hstep\ :\ ∀\ j,\ r\ (j\ +\ 1)\ ≤}
\leanline{6}{(1\ -\ t)\ /\ (1\ +\ t)\ *\ r\ j\ +\ a\ *\ t\ \textasciicircum{}\ 2\ /\ (1\ +\ t))\ (n\ :\ ℕ)\ :}
\leanline{4}{r\ n\ ≤\ Real.exp\ (-2\ *\ t\ *\ (n\ :\ ℝ))\ *\ r\ 0\ +\ a\ *\ t\ /\ 2}
\end{leancode}
\leansource{\texttt{lmc\_\allowbreak{}constant\_\allowbreak{}step} \quad AverageLMC/\allowbreak{}Theorem.lean, lines 84–111}
\begin{leancode}
\leanline{0}{theorem\ lmc\_constant\_step\ \{p\ :\ ℕ\}\ (hp\ :\ 0\ <\ p)\ (f\ :\ State\ p\ →\ ℝ)}
\leanline{4}{(m\ h\ :\ ℝ)\ (L\ :\ ℝ≥0)\ (hm\ :\ 0\ <\ m)\ (hh\ :\ 0\ ≤\ h)}
\leanline{4}{(hf\ :\ Differentiable\ ℝ\ f)\ (hsc\ :\ StrongConvexOn\ univ\ m\ f)}
\leanline{4}{(hLip\ :\ LipschitzWith\ L\ (gradient\ f))\ (hstep\ :\ h\ *\ (L\ :\ ℝ)\ ≤\ 1)}
\leanline{4}{(μ\ :\ Measure\ (State\ p))\ [IsProbabilityMeasure\ μ]}
\leanline{4}{(hfinite\ :\ wasserstein2\ μ\ (gibbsMeasure\ f)\ ≠\ ⊤)\ (k\ :\ ℕ)\ :}
\leanline{4}{wasserstein2Real\ ((ulaLawStep\ f\ h)\textasciicircum{}[k]\ μ)\ (gibbsMeasure\ f)\ \textasciicircum{}\ 2\ ≤}
\leanline{6}{Real.exp\ (-2\ *\ m\ *\ (k\ :\ ℝ)\ *\ h)\ *\ wasserstein2Real\ μ\ (gibbsMeasure\ f)\ \textasciicircum{}\ 2\ +}
\leanline{8}{(averageSmoothness\ f\ +\ m)\ *\ h\ *\ p\ /\ (2\ *\ m)}
\end{leancode}
\leanlabel{How the Lean proof follows the argument}
\texttt{lmc\_\allowbreak{}one\_\allowbreak{}step\_\allowbreak{}stable} first proves that the score correction is nonpositive. \texttt{lmc\_\allowbreak{}constant\_\allowbreak{}step} defines the actual squared-distance sequence as \texttt{r}; \texttt{lmc\allowbreak{}Laws\_\allowbreak{}finite} allows the one-step theorem to be applied at every index. The block \texttt{hs} is the scalar recurrence after normalization. \texttt{constant\_\allowbreak{}step\_\allowbreak{}normalized} iterates it and uses \texttt{contraction\_\allowbreak{}le\_\allowbreak{}exp}. Finally \texttt{lmc\allowbreak{}Laws\_\allowbreak{}const} identifies the recursively constructed marginal laws with iteration of the Gaussian transition.
\end{leanbox}

\section{\texorpdfstring{The decreasing-step schedule}{The decreasing-step schedule}}


Normalize
\(r_k=mW_2^2(\nu_k,\pi)/p\), \(a=\kappa_{\rm av}+1\) and
\(t_k=mh_k=1/(\kappa_\infty+(k-k_0)_+)\). Since \(h_kM_\infty\le1\), the
stable one-step bound reads


\[r_{k+1}\le\frac{1-t_k}{1+t_k}\,r_k+\frac{a\,t_k^2}{1+t_k}.\]


Before \(k_0\) the step is constant, \(t_k=1/\kappa_\infty\le1\), so the
constant-step estimate gives
\(r_{k_0}\le e^{-2k_0/\kappa_\infty}\varepsilon_0^2+a/(2\kappa_\infty)\). The
ceiling defining \(k_0\) makes
\(e^{-2k_0/\kappa_\infty}\varepsilon_0^2\le a/(2\kappa_\infty)\), hence
\(r_{k_0}\le a/\kappa_\infty\). After the warmup, \(r\le a/x\) with \(x\ge1\)
implies


\[\frac{1-1/x}{1+1/x}\,r+\frac{a/x^2}{1+1/x}
\le\frac{x-1}{x+1}\cdot\frac ax+\frac a{x(x+1)}=\frac a{x+1}.\]


Induction with \(x=\kappa_\infty+k-k_0\) gives
\(r_k\le a/(\kappa_\infty+k-k_0)\) for \(k\ge k_0\). Multiplying by \(p/m\)
is Proposition 1. When \(\varepsilon_0=0\), \(k_0=0\) and the induction starts
at once.




\textbf{The warmup ceiling.} For \(\varepsilon_0\ne0\), set \(R=2\kappa_\infty\varepsilon_0^2/a\) and \(\log_+R=\max(\log R,0)\). If \(R\le1\), the initial error already satisfies \(\varepsilon_0^2\le a/(2\kappa_\infty)\), and the warmup may be zero. If \(R>1\), the ceiling ensures \(2k_0/\kappa_\infty\ge\log R\); exponentiating its negative gives \(e^{-2k_0/\kappa_\infty}\varepsilon_0^2\le a/(2\kappa_\infty)\). In both cases the constant-step bias contributes the other half of \(a/\kappa_\infty\). After warmup \(x=\kappa_\infty+k-k_0\ge1\), so the factor \((x-1)/(x+1)\) multiplying the induction hypothesis is nonnegative. The rational calculation above therefore preserves the inequality. Positive dimension and positive \(m\) justify the final multiplication by \(p/m\). When \(\varepsilon_0=0\), the logarithmic warmup is defined to be zero and \(r_0=0\le a/\kappa_\infty\), giving the same induction directly.


\begin{leanbox}{Lean: \texttt{warmup\_\allowbreak{}decay}, \texttt{harmonic\_\allowbreak{}step}, \texttt{lmc\_\allowbreak{}variable\_\allowbreak{}step}}
\leansource{\texttt{warmup\_\allowbreak{}decay} \quad AverageLMC/\allowbreak{}Complexity.lean, lines 39–56}
\begin{leancode}
\leanline{0}{theorem\ warmup\_decay\ \{κ\ a\ ε₀\ :\ ℝ\}\ (hκ\ :\ 0\ <\ κ)\ (ha\ :\ 0\ <\ a)}
\leanline{4}{(hε₀\ :\ 0\ <\ ε₀)\ :}
\leanline{4}{Real.exp\ (-2\ *\ (warmup\ κ\ a\ ε₀\ :\ ℝ)\ /\ κ)\ *\ ε₀\ \textasciicircum{}\ 2\ ≤\ a\ /\ (2\ *\ κ)}
\end{leancode}
\leansource{\texttt{harmonic\_\allowbreak{}step} \quad AverageLMC/\allowbreak{}Recurrence.lean, lines 86–100}
\begin{leancode}
\leanline{0}{theorem\ harmonic\_step\ \{a\ x\ r\ :\ ℝ\}\ (hx\ :\ 1\ ≤\ x)}
\leanline{4}{(hr\ :\ r\ ≤\ a\ /\ x)\ :}
\leanline{4}{(1\ -\ 1\ /\ x)\ /\ (1\ +\ 1\ /\ x)\ *\ r\ +\ a\ *\ (1\ /\ x)\ \textasciicircum{}\ 2\ /\ (1\ +\ 1\ /\ x)}
\leanline{6}{≤\ a\ /\ (x\ +\ 1)}
\end{leancode}
\leansource{\texttt{lmc\_\allowbreak{}variable\_\allowbreak{}step} \quad AverageLMC/\allowbreak{}Theorem.lean, lines 168–227}
\begin{leancode}
\leanline{0}{theorem\ lmc\_variable\_step\ \{p\ :\ ℕ\}\ (hp\ :\ 0\ <\ p)\ (f\ :\ State\ p\ →\ ℝ)}
\leanline{4}{(m\ ε₀\ :\ ℝ)\ (L\ :\ ℝ≥0)\ (hm\ :\ 0\ <\ m)}
\leanline{4}{(hf\ :\ Differentiable\ ℝ\ f)\ (hsc\ :\ StrongConvexOn\ univ\ m\ f)}
\leanline{4}{(hLip\ :\ LipschitzWith\ L\ (gradient\ f))}
\leanline{4}{(μ\ :\ Measure\ (State\ p))\ [IsProbabilityMeasure\ μ]}
\leanline{4}{(hfinite\ :\ wasserstein2\ μ\ (gibbsMeasure\ f)\ ≠\ ⊤)}
\leanline{4}{(hinit\ :\ wasserstein2Real\ μ\ (gibbsMeasure\ f)\ \textasciicircum{}\ 2\ ≤\ ε₀\ \textasciicircum{}\ 2\ *\ p\ /\ m)}
\leanline{4}{(k\ :\ ℕ)\ (hk\ :\ warmup\ ((L\ :\ ℝ)\ /\ m)\ (averageSmoothness\ f\ /\ m\ +\ 1)\ ε₀\ ≤\ k)\ :}
\leanline{4}{let\ κ\ :=\ (L\ :\ ℝ)\ /\ m}
\leanline{4}{let\ a\ :=\ averageSmoothness\ f\ /\ m\ +\ 1}
\leanline{4}{let\ k₀\ :=\ warmup\ κ\ a\ ε₀}
\leanline{4}{wasserstein2Real\ (lmcLaws\ f\ (lmcVariableSchedule\ (L\ :\ ℝ)\ m\ k₀)\ μ\ k)}
\leanline{6}{(gibbsMeasure\ f)\ \textasciicircum{}\ 2\ ≤}
\leanline{6}{a\ /\ (κ\ +\ ((k\ -\ k₀\ :\ ℕ)\ :\ ℝ))\ *\ p\ /\ m}
\end{leancode}
\leanlabel{How the Lean proof follows the argument}
\texttt{warmup\_\allowbreak{}decay} checks the initial exponential allocation, including the logarithmic ceiling. \texttt{harmonic\_\allowbreak{}step} proves the single rational induction step after warmup. \texttt{lmc\_\allowbreak{}variable\_\allowbreak{}step} verifies admissibility of every actual step in \texttt{hs}, normalizes the initial cost in \texttt{hi}, and derives the normalized recurrence in \texttt{hrec}. \texttt{variable\_\allowbreak{}step\_\allowbreak{}normalized\_\allowbreak{}all} handles zero initial error explicitly before applying the scalar warmup-and-tail estimate.
\end{leanbox}

\section{\texorpdfstring{The Laplacian refinement}{The Laplacian refinement}}


Use the coupling and
stationary diffusion of the first step, with \(L^2\) inner products. Put
\(\Psi=(\vartheta-h\nabla f(\vartheta))-(Y_0-h\nabla f(Y_0))\) and
\(R=\int_0^h(\nabla f(Y_t)-\nabla f(Y_0))\,dt\). The output error is
\(A+B_h-Y_h=\Psi+R\), and for \(hM_\infty\le1\) the gradient step contracts:
\(\|\Psi\|\le(1-mh)\|\vartheta-Y_0\|\).




The diffusion generator maps \(\nabla f\) to
\(\nabla\Delta f-\nabla^2f\,\nabla f\). Because \(\Psi\) is independent of the
Brownian motion, a weighted Dynkin identity, stationarity,
\(\|\nabla\Delta f\|\le M_\Delta\sqrt p\) and the Gibbs score moment
\(\mathbb E_\pi\|\nabla f\|^2\le pM_{\rm av}\) give


\[\langle\Psi,\nabla f(Y_t)-\nabla f(Y_0)\rangle\le t\,\|\Psi\|\,G_\Delta,
\qquad G_\Delta=M_\infty\sqrt{M_{\rm av}p}+M_\Delta\sqrt p.\]


The Dynkin identity is derived for compactly supported \(C^2\) tests from
fresh Gaussian Euler steps. Cutoffs of the score components and radial
clipping of \(\Psi\) extend it to this pairing. The integral equation,
Brownian moments and \(\mathbb E_\pi\|\nabla f\|^2\le pM_\infty\) give
\(\|R\|\le\tfrac32M_\infty h^{3/2}\sqrt p\). Expanding \(\|\Psi+R\|^2\) and
taking the coupling infimum, \(e_k=W_2(\nu_k,\pi)\) and \(q=1-mh\) satisfy


\[e_{k+1}^2\le q^2e_k^2+q\,h^2G_\Delta\,e_k+\bigl(\tfrac32M_\infty h^{3/2}\sqrt p\bigr)^2.\]


A recurrence \(e_{k+1}^2\le q^2e_k^2+2qCe_k+B^2\) with
\(C\le(1-q)a\) and \(B^2\le(1-q^2)b^2\) gives \(e_k\le q^ke_0+a+b\). Take
\(a=hG_\Delta/(2m)\) and \(b=3hM_\infty\sqrt{M_{\rm av}p}/(2m)\); the second
condition uses \(m\le M_{\rm av}\). Then


\[a+b=\frac{(4M_\infty\sqrt{M_{\rm av}}+M_\Delta)h\sqrt p}{2m}
\le\frac{(2M_\infty\sqrt{M_{\rm av}}+M_\Delta)h\sqrt p}{m},\]


which is Theorem 2 for \(0<h<1/M_\infty\).




\textbf{Bounds used by the recurrence.} Here the additional assumptions are \(f\in C^3\), \(M_\Delta\ge0\) and \(\|\nabla\Delta f(x)\|^2\le M_\Delta^2p\) for every \(x\), with \(0<hM_\infty<1\). The derivative of the gradient has operator norm at most \(M_\infty\). Thus Cauchy–Schwarz and the score moment bound give

\[\|\nabla\Delta f(Y_t)-\nabla^2f(Y_t)\nabla f(Y_t)\|_{L^2}\le(M_\Delta+M_\infty\sqrt{M_{\rm av}})\sqrt p=G_\Delta.\]

Integrating the pairing bound over time gives \(2\langle\Psi,R\rangle\le h^2G_\Delta\|\Psi\|_{L^2}\). For the total defect, the integral equation, Minkowski and the stationary score moment give \(\|Y_t-Y_0\|_{L^2}\le t\sqrt{M_\infty p}+\sqrt{2pt}\). Hence

\[\begin{aligned}\|R\|_{L^2}&\le M_\infty\int_0^h\left(t\sqrt{M_\infty p}+\sqrt{2pt}\right)dt\\&=M_\infty\left(\frac{h^2}2\sqrt{M_\infty p}+\frac{2\sqrt2}3h^{3/2}\sqrt p\right)\\&\le\left(\frac12+\frac{2\sqrt2}3\right)M_\infty h^{3/2}\sqrt p\le\frac32M_\infty h^{3/2}\sqrt p.\end{aligned}\]

The penultimate inequality uses \(hM_\infty\le1\); the last uses \(\sqrt2\le3/2\). For the stated recurrence floor, \(C=h^2G_\Delta/2=(1-q)a\). Also \(B^2\le mh\,b^2\) follows from \(m\le M_{\rm av}\), and \(mh\le1-q^2=mh(2-mh)\). These verify both scalar recurrence hypotheses. The final enlargement of the nonnegative \(M_\Delta\) term gives precisely the printed coefficient.


\begin{leanbox}{Lean: \texttt{lp\_\allowbreak{}stationary\_\allowbreak{}gradient\_\allowbreak{}pairing\_\allowbreak{}bound}, \texttt{lp\_\allowbreak{}laplacian\_\allowbreak{}one\_\allowbreak{}step\_\allowbreak{}bound}, \texttt{laplacian\_\allowbreak{}recurrence\_\allowbreak{}bound\_\allowbreak{}of\_\allowbreak{}global\_\allowbreak{}defect}}
\leansource{\texttt{lp\_\allowbreak{}stationary\_\allowbreak{}gradient\_\allowbreak{}pairing\_\allowbreak{}bound} \quad AverageLMC/\allowbreak{}LaplacianPairing.lean, lines 50–210}
\begin{leancode}
\leanline{0}{theorem\ lp\_stationary\_gradient\_pairing\_bound\ \{p\ :\ ℕ\}\ (hp\ :\ 0\ <\ p)}
\leanline{4}{(f\ :\ State\ p\ →\ ℝ)\ (m\ :\ ℝ)\ (L\ :\ ℝ≥0)\ (hm\ :\ 0\ <\ m)}
\leanline{4}{(hsc\ :\ StrongConvexOn\ univ\ m\ f)\ (hLip\ :\ LipschitzWith\ L\ (gradient\ f))}
\leanline{4}{(Mdelta\ :\ ℝ)\ (hMdelta\ :\ 0\ ≤\ Mdelta)\ (hδ\ :\ LaplacianSmooth\ f\ Mdelta)}
\leanline{4}{(initial\ :\ Lp\ (State\ p)\ 2\ P)}
\leanline{4}{(B\ :\ Fin\ p\ →\ ℝ≥0\ →\ Ω\ →\ ℝ)\ (hB\ :\ ∀\ i,\ IsPreBrownianReal\ (B\ i)\ P)}
\leanline{4}{(hscalar\ :\ ∀\ r\ :\ ℝ,\ 0\ <\ r\ →\ iIndepFun\ (fun\ q\ :\ ℕ\ ×\ Fin\ p\ ↦\ fun\ a\ ↦}
\leanline{6}{ULA.Hilbert.brownianEulerNoise\ P\ B\ hB\ r\ q.1\ a\ q.2)\ P)}
\leanline{4}{(Ψ\ :\ Lp\ (State\ p)\ 2\ P)}
\leanline{4}{(hind\ :\ ∀\ r\ :\ ℝ,\ 0\ <\ r\ →\ IndepFun\ (fun\ a\ ↦\ (Ψ\ a,\ initial\ a))}
\leanline{6}{(fun\ a\ k\ ↦\ ULA.Hilbert.brownianEulerNoise\ P\ B\ hB\ r\ k\ a)\ P)}
\leanline{4}{(Y\ :\ ℝ\ →\ Lp\ (State\ p)\ 2\ P)\ (hY\ :\ Continuous\ Y)}
\leanline{4}{(T\ :\ ℝ)\ (hT\ :\ 0\ ≤\ T)}
\leanline{4}{(hYe\ :\ ∀\ t\ ∈\ Icc\ 0\ T,\ Y\ t\ -\ initial\ =}
\leanline{6}{-(∫\ s\ in\ 0..t,\ ULA.lpDrift\ P\ (gradient\ f)\ hLip\ (Y\ s))\ +}
\leanline{8}{ULA.brownianLangevinNoise\ P\ B\ hB\ 0\ t)}
\leanline{4}{(hstationary\ :\ ∀\ t\ ∈\ Icc\ 0\ T,\ P.map\ (Y\ t)\ =\ ULA.gibbsMeasure\ f)\ :}
\leanline{4}{inner\ ℝ\ Ψ\ (ULA.lpDrift\ P\ (gradient\ f)\ hLip\ (Y\ T)\ -}
\leanline{6}{ULA.lpDrift\ P\ (gradient\ f)\ hLip\ initial)\ ≤}
\leanline{8}{T\ *\ ‖Ψ‖\ *\ ((L\ :\ ℝ)\ *\ Real.sqrt\ (averageSmoothness\ f\ *\ (p\ :\ ℝ))\ +}
\leanline{10}{Mdelta\ *\ Real.sqrt\ (p\ :\ ℝ))}
\end{leancode}
\leansource{\texttt{lp\_\allowbreak{}laplacian\_\allowbreak{}one\_\allowbreak{}step\_\allowbreak{}bound} \quad AverageLMC/\allowbreak{}LaplacianOneStep.lean, lines 21–88}
\begin{leancode}
\leanline{0}{theorem\ lp\_laplacian\_one\_step\_bound}
\leanline{4}{(f\ :\ E\ →\ ℝ)\ (m\ h\ d\ Rgen\ :\ ℝ)\ (M\ :\ ℝ≥0)}
\leanline{4}{(hf\ :\ Differentiable\ ℝ\ f)\ (hsc\ :\ StrongConvexOn\ univ\ m\ f)}
\leanline{4}{(hLip\ :\ LipschitzWith\ M\ (gradient\ f))\ (hm\ :\ 0\ <\ m)\ (hmM\ :\ m\ ≤\ (M\ :\ ℝ))}
\leanline{4}{(hh\ :\ 0\ ≤\ h)\ (hd\ :\ 0\ ≤\ d)\ (hsmall\ :\ h\ *\ (M\ :\ ℝ)\ ≤\ 1)\ (hRgen\ :\ 0\ ≤\ Rgen)}
\leanline{4}{(X\ :\ Lp\ E\ 2\ P)\ (Y\ B\ :\ ℝ\ →\ Lp\ E\ 2\ P)\ (hY\ :\ Continuous\ Y)}
\leanline{4}{(hG\ :\ ∀\ t\ ∈\ Icc\ 0\ h,\ ‖ULA.lpDrift\ P\ (gradient\ f)\ hLip\ (Y\ t)‖\ ≤}
\leanline{6}{Real.sqrt\ ((M\ :\ ℝ)\ *\ d))}
\leanline{4}{(hB\ :\ ∀\ t\ ∈\ Icc\ 0\ h,\ ‖B\ t‖\ ≤\ Real.sqrt\ (2\ *\ d)\ *\ Real.sqrt\ t)}
\leanline{4}{(hYe\ :\ ∀\ t\ ∈\ Icc\ 0\ h,\ Y\ t\ -\ Y\ 0\ =}
\leanline{6}{-(∫\ s\ in\ 0..t,\ ULA.lpDrift\ P\ (gradient\ f)\ hLip\ (Y\ s))\ +\ B\ t)}
\leanline{4}{(hpair\ :\ ∀\ t\ ∈\ Icc\ 0\ h,}
\leanline{6}{{\leanfallback ⟪}(X\ -\ h\ •\ ULA.lpDrift\ P\ (gradient\ f)\ hLip\ X)\ -}
\leanline{10}{(Y\ 0\ -\ h\ •\ ULA.lpDrift\ P\ (gradient\ f)\ hLip\ (Y\ 0)),}
\leanline{8}{ULA.lpDrift\ P\ (gradient\ f)\ hLip\ (Y\ t)\ -}
\leanline{10}{ULA.lpDrift\ P\ (gradient\ f)\ hLip\ (Y\ 0){\leanfallback ⟫}\ ≤}
\leanline{8}{t\ *\ ‖(X\ -\ h\ •\ ULA.lpDrift\ P\ (gradient\ f)\ hLip\ X)\ -}
\leanline{10}{(Y\ 0\ -\ h\ •\ ULA.lpDrift\ P\ (gradient\ f)\ hLip\ (Y\ 0))‖\ *\ Rgen)\ :}
\leanline{4}{‖X\ -\ h\ •\ ULA.lpDrift\ P\ (gradient\ f)\ hLip\ X\ +\ B\ h\ -\ Y\ h‖\ \textasciicircum{}\ 2\ ≤}
\leanline{6}{(1\ -\ m\ *\ h)\ \textasciicircum{}\ 2\ *\ ‖X\ -\ Y\ 0‖\ \textasciicircum{}\ 2\ +}
\leanline{6}{h\ \textasciicircum{}\ 2\ *\ (1\ -\ m\ *\ h)\ *\ ‖X\ -\ Y\ 0‖\ *\ Rgen\ +}
\leanline{6}{((3\ /\ 2\ :\ ℝ)\ *\ (M\ :\ ℝ)\ *\ h\ *\ Real.sqrt\ h\ *\ Real.sqrt\ d)\ \textasciicircum{}\ 2}
\end{leancode}
\leansource{\texttt{laplacian\_\allowbreak{}recurrence\_\allowbreak{}bound\_\allowbreak{}of\_\allowbreak{}global\_\allowbreak{}defect} \quad AverageLMC/\allowbreak{}LaplacianLocalError.lean, lines 57–116}
\begin{leancode}
\leanline{0}{theorem\ laplacian\_recurrence\_bound\_of\_global\_defect}
\leanline{4}{(e\ :\ ℕ\ →\ ℝ)\ (m\ M\ h\ A\ δ\ S\ :\ ℝ)}
\leanline{4}{(he\ :\ ∀\ k,\ 0\ ≤\ e\ k)\ (hm\ :\ 0\ <\ m)\ (hmM\ :\ m\ ≤\ M)}
\leanline{4}{(hh\ :\ 0\ ≤\ h)\ (hhM\ :\ h\ *\ M\ ≤\ 1)}
\leanline{4}{(hA\ :\ 0\ ≤\ A)\ (hmA\ :\ m\ ≤\ A\ \textasciicircum{}\ 2)\ (hδ\ :\ 0\ ≤\ δ)\ (hS\ :\ 0\ ≤\ S)}
\leanline{4}{(hstep\ :\ ∀\ k,\ (e\ (k\ +\ 1))\ \textasciicircum{}\ 2\ ≤}
\leanline{6}{(1\ -\ m\ *\ h)\ \textasciicircum{}\ 2\ *\ (e\ k)\ \textasciicircum{}\ 2\ +}
\leanline{8}{(1\ -\ m\ *\ h)\ *\ h\ \textasciicircum{}\ 2\ *\ (M\ *\ A\ +\ δ)\ *\ S\ *\ e\ k\ +}
\leanline{10}{((3\ /\ 2\ :\ ℝ)\ *\ M\ *\ h\ *\ Real.sqrt\ h\ *\ S)\ \textasciicircum{}\ 2)\ (n\ :\ ℕ)\ :}
\leanline{4}{e\ n\ ≤\ (1\ -\ m\ *\ h)\ \textasciicircum{}\ n\ *\ e\ 0\ +\ (2\ *\ M\ *\ A\ +\ δ)\ *\ h\ *\ S\ /\ m}
\end{leancode}
\leanlabel{How the Lean proof follows the argument}
\texttt{lp\_\allowbreak{}stationary\_\allowbreak{}gradient\_\allowbreak{}pairing\_\allowbreak{}bound} bounds the cross term with the actual stationary diffusion and the stated gradient-of-Laplacian assumption. \texttt{lp\_\allowbreak{}laplacian\_\allowbreak{}one\_\allowbreak{}step\_\allowbreak{}bound} combines that pairing with the local integral remainder and gradient-step contraction. \texttt{laplacian\_\allowbreak{}recurrence\_\allowbreak{}bound\_\allowbreak{}of\_\allowbreak{}global\_\allowbreak{}defect} iterates the resulting quadratic-plus-linear scalar estimate. The final theorem \texttt{lmc\_\allowbreak{}laplacian} obtains the required finite-distance recurrence for the actual Gaussian laws and simplifies its constant using the exact Gibbs score moment bound.
\end{leanbox}

\section{\texorpdfstring{Control-variate SGLD}{Control-variate SGLD}}


Let
\(g_I=n(\nabla f_I(\vartheta)-\nabla f_I(\theta_*))\). Since
\(\nabla f(\theta_*)=0\), its mean over the uniform index is
\(\nabla f(\vartheta)\). The SGLD output is the LMC output plus
\(-h(g_I-\nabla f(\vartheta))\), which has conditional mean zero given
\((\vartheta,\xi)\). Lifting any coupling of the LMC output with \(\pi\) to the
index by disintegration, the cross term vanishes, so


\[W_2^2(\nu^+,\pi)\le W_2^2(\nu^+_{\rm LMC},\pi)
+h^2\bigl(\mathbb E\|g_I\|^2-\mathbb E\|\nabla f(\vartheta)\|^2\bigr).\]


Write \(m=n\bar m\) and let \(M\) be a Lipschitz constant of \(\nabla f\).
The one-step LMC bound turns the score coefficient into
\(h^2(hM-1)/(1+mh)-h^2\le h^2(hM-2)/(1+mh)\), which is Proposition 3's first
claim. The component seminorms give
\(\mathbb E_I\|g_I\|^2\le n^2M_{\rm rms}^2\|\vartheta-\theta_*\|^2\), and the
Gibbs bound \(\mathbb E_\pi\|Y-\theta_*\|^2\le p/m\) gives
\(\mathbb E\|\vartheta-\theta_*\|^2\le2W_2^2(\nu,\pi)+2p/m\). Take
\(M=R=nM_{\rm rms}\). Because \(m\le R\), the step condition
\(3R^2h\le m\) implies \(hR\le1/3\), so the score term is nonpositive and


\[r_{k+1}\le\Bigl(\frac{1-mh}{1+mh}+2R^2h^2\Bigr)r_k
+\Bigl(\frac{M_{\rm av}(f)+m}{1+mh}+\frac{2R^2}m\Bigr)h^2p.\]


The step condition makes the contraction at most \((1+mh)^{-1}\) and the
forcing at most \((M_{\rm av}(f)+2m+2R^2/m)h^2p/(1+mh)\). The fixed point is
\((M_{\rm av}(f)+2m+2R^2/m)hp/m\). Substituting \(m=n\bar m\),
\(R=nM_{\rm rms}\) and \(M_{\rm av}(f)=n\bar M_{\rm av}\) gives the SGLD
bound; an infinite initial distance makes it trivial.




\textbf{Why the variance adds.} Conditional on \((\vartheta,\xi)\), the index \(I\) is uniform and its centered error \(\eta=g_I-\nabla f(\vartheta)\) has mean zero. Expanding a lifted coupled square gives \(\mathbb E\|D-h\eta\|^2=\mathbb E\|D\|^2+h^2\mathbb E\|\eta\|^2\), because the conditional expectation of \(\langle D,\eta\rangle\) is zero. Also \(\mathbb E_I\|\eta\|^2=\mathbb E_I\|g_I\|^2-\|\nabla f(\vartheta)\|^2\). If \(L_i\) is the actual Lipschitz seminorm of \(\nabla f_i\), then

\[\mathbb E_I\|g_I\|^2=n\sum_i\|\nabla f_i(\vartheta)-\nabla f_i(\theta_*)\|^2\le n\sum_iL_i^2\|\vartheta-\theta_*\|^2=n^2M_{\rm rms}^2\|\vartheta-\theta_*\|^2.\]

For any coupling with Gibbs, \(\|\vartheta-\theta_*\|^2\le2\|\vartheta-Y\|^2+2\|Y-\theta_*\|^2\). Integrate, use the Gibbs second moment, and pass costs to their infimum to obtain the displayed second-moment bound. With \(R=nM_{\rm rms}\), the small-step condition \(3R^2h\le m\) and \(m\le R\) imply \(mh\le1/3\). The forcing comparison uses \(2R^2h\le m\): after multiplication by \(1+mh\), its extra term is at most \(mh\) in the required coefficient. The recurrence fixed point is the forcing divided by \(1-(1+mh)^{-1}=mh/(1+mh)\). At \(h=0\) the transition is the identity and this division is unnecessary.


\begin{leanbox}{Lean: \texttt{sgld\_\allowbreak{}variance\_\allowbreak{}addition}, \texttt{controlVariateEstimator\_\allowbreak{}second\_\allowbreak{}moment\_\allowbreak{}le}, \texttt{sgld\_\allowbreak{}one\_\allowbreak{}step\_\allowbreak{}rms} and 1 more}
\leansource{\texttt{sgld\_\allowbreak{}variance\_\allowbreak{}addition} \quad AverageLMC/\allowbreak{}SGLDLaw.lean, lines 98–141}
\begin{leancode}
\leanline{0}{theorem\ sgld\_variance\_addition\ \{p\ n\ :\ ℕ\}\ (hn\ :\ 0\ <\ n)}
\leanline{4}{(f\ :\ Fin\ n\ →\ State\ p\ →\ ℝ)\ (hf\ :\ ∀\ i,\ Differentiable\ ℝ\ (f\ i))}
\leanline{4}{(L\ :\ Fin\ n\ →\ ℝ≥0)\ (hLip\ :\ ∀\ i,\ LipschitzWith\ (L\ i)\ (gradient\ (f\ i)))}
\leanline{4}{(M\ :\ ℝ≥0)\ (hfull\ :\ LipschitzWith\ M\ (gradient\ (componentPotential\ f)))}
\leanline{4}{(anchor\ :\ State\ p)\ (hstar\ :\ gradient\ (componentPotential\ f)\ anchor\ =\ 0)}
\leanline{4}{(h\ :\ ℝ)\ (μ\ π\ :\ Measure\ (State\ p))\ [IsProbabilityMeasure\ μ]\ [IsProbabilityMeasure\ π]}
\leanline{4}{(hμ\ :\ MemLp\ (fun\ x\ :\ State\ p\ ↦\ x)\ 2\ μ)\ (hπ\ :\ MemLp\ (fun\ x\ :\ State\ p\ ↦\ x)\ 2\ π)\ :}
\leanline{4}{ULA.wasserstein2Real\ (sgldLawStep\ f\ anchor\ h\ μ)\ π\ \textasciicircum{}\ 2\ ≤}
\leanline{6}{ULA.wasserstein2Real\ (ULA.ulaLawStep\ (componentPotential\ f)\ h\ μ)\ π\ \textasciicircum{}\ 2\ +}
\leanline{8}{h\ \textasciicircum{}\ 2\ *\ ∫\ q\ :\ (State\ p\ ×\ State\ p)\ ×\ Fin\ n,}
\leanline{10}{‖controlVariateError\ f\ anchor\ q.1.1\ q.2‖\ \textasciicircum{}\ 2}
\leanline{12}{∂(μ.prod\ (stdGaussian\ (State\ p))).prod\ (uniformIndexLaw\ n)}
\end{leancode}
\leansource{\texttt{controlVariateEstimator\_\allowbreak{}second\_\allowbreak{}moment\_\allowbreak{}le} \quad AverageLMC/\allowbreak{}SGLDEstimator.lean, lines 152–167}
\begin{leancode}
\leanline{0}{theorem\ controlVariateEstimator\_second\_moment\_le\ \{p\ n\ :\ ℕ\}\ (hp\ :\ 0\ <\ p)\ (hn\ :\ 0\ <\ n)}
\leanline{4}{(f\ :\ Fin\ n\ →\ State\ p\ →\ ℝ)\ (L\ :\ Fin\ n\ →\ ℝ≥0)}
\leanline{4}{(hLip\ :\ ∀\ i,\ LipschitzWith\ (L\ i)\ (gradient\ (f\ i)))\ (xstar\ x\ :\ State\ p)\ :}
\leanline{4}{(∫\ i,\ ‖controlVariateEstimator\ f\ xstar\ x\ i‖\ \textasciicircum{}\ 2\ ∂uniformIndexLaw\ n)\ ≤}
\leanline{6}{(n\ :\ ℝ)\ \textasciicircum{}\ 2\ *\ rmsComponentSmoothness\ f\ \textasciicircum{}\ 2\ *\ ‖x\ -\ xstar‖\ \textasciicircum{}\ 2}
\end{leancode}
\leansource{\texttt{sgld\_\allowbreak{}one\_\allowbreak{}step\_\allowbreak{}rms} \quad AverageLMC/\allowbreak{}SGLDTheorem.lean, lines 56–103}
\begin{leancode}
\leanline{0}{theorem\ sgld\_one\_step\_rms\ \{p\ n\ :\ ℕ\}\ (hp\ :\ 0\ <\ p)\ (hn\ :\ 0\ <\ n)}
\leanline{4}{(f\ :\ Fin\ n\ →\ State\ p\ →\ ℝ)\ (hf\ :\ ∀\ i,\ Differentiable\ ℝ\ (f\ i))}
\leanline{4}{(L\ :\ Fin\ n\ →\ ℝ≥0)\ (hLip\ :\ ∀\ i,\ LipschitzWith\ (L\ i)\ (gradient\ (f\ i)))}
\leanline{4}{(m\ h\ :\ ℝ)\ (M\ :\ ℝ≥0)\ (hm\ :\ 0\ <\ m)\ (hh\ :\ 0\ ≤\ h)}
\leanline{4}{(hsc\ :\ StrongConvexOn\ univ\ m\ (componentPotential\ f))}
\leanline{4}{(hfull\ :\ LipschitzWith\ M\ (gradient\ (componentPotential\ f)))}
\leanline{4}{(hstep\ :\ h\ *\ (M\ :\ ℝ)\ ≤\ 2)}
\leanline{4}{(anchor\ :\ State\ p)\ (hanchor\ :\ ∀\ x,\ componentPotential\ f\ anchor\ ≤\ componentPotential\ f\ x)}
\leanline{4}{(hstar\ :\ gradient\ (componentPotential\ f)\ anchor\ =\ 0)}
\leanline{4}{(μ\ :\ Measure\ (State\ p))\ [IsProbabilityMeasure\ μ]}
\leanline{4}{(hfinite\ :\ wasserstein2\ μ\ (gibbsMeasure\ (componentPotential\ f))\ ≠\ ⊤)\ :}
\leanline{4}{let\ R\ :=\ (n\ :\ ℝ)\ *\ rmsComponentSmoothness\ f}
\leanline{4}{wasserstein2Real\ (sgldLawStep\ f\ anchor\ h\ μ)\ (gibbsMeasure\ (componentPotential\ f))\ \textasciicircum{}\ 2\ ≤}
\leanline{6}{((1\ -\ m\ *\ h)\ /\ (1\ +\ m\ *\ h)\ +\ 2\ *\ R\ \textasciicircum{}\ 2\ *\ h\ \textasciicircum{}\ 2)\ *}
\leanline{8}{wasserstein2Real\ μ\ (gibbsMeasure\ (componentPotential\ f))\ \textasciicircum{}\ 2\ +}
\leanline{6}{((averageSmoothness\ (componentPotential\ f)\ +\ m)\ /\ (1\ +\ m\ *\ h)\ +\ 2\ *\ R\ \textasciicircum{}\ 2\ /\ m)\ *}
\leanline{8}{h\ \textasciicircum{}\ 2\ *\ p}
\end{leancode}
\leansource{\texttt{sgld\_\allowbreak{}recurrence\_\allowbreak{}bound} \quad AverageLMC/\allowbreak{}SGLDRecurrence.lean, lines 52–80}
\begin{leancode}
\leanline{0}{theorem\ sgld\_recurrence\_bound\ (r\ :\ ℕ\ →\ ℝ)\ (m\ R\ A\ h\ p\ :\ ℝ)}
\leanline{4}{(hm\ :\ 0\ <\ m)\ (hh\ :\ 0\ ≤\ h)\ (hp\ :\ 0\ ≤\ p)\ (hA\ :\ 0\ ≤\ A)}
\leanline{4}{(hR\ :\ m\ \textasciicircum{}\ 2\ ≤\ R\ \textasciicircum{}\ 2)\ (hstep\ :\ 3\ *\ R\ \textasciicircum{}\ 2\ *\ h\ ≤\ m)}
\leanline{4}{(hr\ :\ ∀\ n,\ 0\ ≤\ r\ n)}
\leanline{4}{(hrec\ :\ ∀\ n,\ r\ (n\ +\ 1)\ ≤}
\leanline{6}{((1\ -\ m\ *\ h)\ /\ (1\ +\ m\ *\ h)\ +\ 2\ *\ R\ \textasciicircum{}\ 2\ *\ h\ \textasciicircum{}\ 2)\ *\ r\ n\ +}
\leanline{6}{((A\ +\ m)\ /\ (1\ +\ m\ *\ h)\ +\ 2\ *\ R\ \textasciicircum{}\ 2\ /\ m)\ *\ h\ \textasciicircum{}\ 2\ *\ p)}
\leanline{4}{(k\ :\ ℕ)\ :}
\leanline{4}{r\ k\ ≤\ (1\ +\ m\ *\ h)⁻¹\ \textasciicircum{}\ k\ *\ r\ 0\ +}
\leanline{6}{(A\ +\ 2\ *\ m\ +\ 2\ *\ R\ \textasciicircum{}\ 2\ /\ m)\ *\ h\ *\ p\ /\ m}
\end{leancode}
\leanlabel{How the Lean proof follows the argument}
\texttt{sgld\_\allowbreak{}variance\_\allowbreak{}addition} uses a lifted coupling to remove the centered estimator cross term. \texttt{control\allowbreak{}Variate\allowbreak{}Estimator\_\allowbreak{}second\_\allowbreak{}moment\_\allowbreak{}le} bounds the estimator using the actual component Lipschitz seminorms. \texttt{sgld\_\allowbreak{}one\_\allowbreak{}step\_\allowbreak{}rms} combines this variance bound with the one-step LMC comparison. \texttt{sgld\_\allowbreak{}recurrence\_\allowbreak{}bound} uses the small-step condition to bound the contraction and forcing, then sums the affine recurrence. The extended theorem separates infinite initial transport before converting any distance to a real number.
\end{leanbox}

\end{document}
